\documentclass{standalone}
\usepackage{standalone}

\begin{document}
\chapter{Introduction}

\section{Project Overview}
Sarabo is a web-based platform for managing electronics repair services, which I developed on my own as my final year project. It handles the whole life of a repair request, from the moment a customer describes a faulty device to the moment the technician gets paid, as one workflow that can be traced from start to finish. Three people are involved: a \emph{Customer} who needs something repaired, a \emph{Technician} who does the repair, and an \emph{Admin} who keeps the service running.

The most important design decision was to treat a repair request as a \emph{case that moves through states}, not as a simple database record. Every important activity moves the case from one clearly defined state to the next: creating the request, technician offers, booking, pickup, inspection, quotation, payment, repair, hand-over, settlement and feedback. For each move, the server checks that this user is allowed to make it at this moment. This is the main difference between Sarabo and a basic create-read-update-delete (CRUD) application. Whether a later action is correct depends on earlier actions having happened, and on them having happened only once.

The project was built in two main stages. In the first stage (July--August 2026) I built the core repair lifecycle, in which an admin assigned a technician to each request. An earlier version of this report described that stage. In the second stage (September--October 2026) Sarabo became a two-sided marketplace. It gained an open jobs board where technicians send offers, a paid inspection fee with automatic refunds, pickup scheduling, customer reviews, email notifications, operational reports, four AI-assisted features and a production deployment. A final phase (11--16 October 2026) then added offer ranking, a repair warranty, deadlines for every waiting state, a minimum inspection fee and faster hosting. This report describes the system as it stands after all of that work.

\section{Problem Statement}
A repair service needs both technical work and a lot of coordination. The customer has to describe the device and the problem, get an idea of the cost before handing the device over, choose someone they can trust, pay safely and find out whether the device has actually been repaired. The technician needs a steady supply of nearby jobs that match their skills, a clear way to record an inspection and a quote, and a reliable way to be paid. The administrator needs to see which requests are stuck, which payments need attention, and whether customers are being turned away.

When all of this is handled through phone calls, messaging apps and paper notes, things easily go wrong. A customer is charged before agreeing to a price. A technician is booked twice for the same time. An inspection fee is neither refunded nor taken off the bill. A technician's earning is paid out twice, or private photos of a customer's device end up in the wrong hands. Sarabo deals with this by modelling the repair process as a workflow that knows each user's role, with clear rules for every change of state, validation on the server and updates that run as transactions.

\section{Motivation}
I had five main reasons for choosing this project:
\begin{enumerate}
  \item to work on a real, local problem, by making electronics repair in Bangladesh more transparent and trustworthy for customers and easier to access for independent technicians;
  \item to build a realistic full-stack system that goes beyond simple CRUD, where roles and state changes are enforced on the server instead of being left to the browser;
  \item to integrate real, production-grade services (identity, card payments and refunds, private file storage, email and generative AI) without exposing any privileged credentials;
  \item to use generative AI where it actually helps people, while keeping it safe: it must not invent prices or leak personal data, and its usage must have hard limits; and
  \item to give every user a clear interface in which the next useful action is always easy to see.
\end{enumerate}

\section{Objectives}
The main objective was to build a reliable platform that manages the complete lifecycle of an electronics repair. Table~\ref{tab:objectives} lists the specific objectives.

\begin{table}[htbp]
\centering
\caption{Project objectives}
\label{tab:objectives}
\begin{tabularx}{\textwidth}{@{}lX@{}}
\toprule
\textbf{No.} & \textbf{Objective} \\
\midrule
O1 & Provide secure authentication and role-aware access for Customer, Technician and Admin. \\
O2 & Let customers create structured repair requests (device, service, region, pickup window, optional damage photos), with optional AI assistance. \\
O3 & Match requests only to approved technicians with the right expertise in the customer's region, through an open offer-based job portal. \\
O4 & Support a paid inspection fee that books a technician, is credited toward the final price, and is refunded or split fairly when the repair does not go ahead. \\
O5 & Provide an inspection and quotation workflow with customer approval, rejection and revision. \\
O6 & Accept payment only for approved amounts, with a single authoritative paid state and downloadable receipts. \\
O7 & Support repair progress, completion evidence, receipt confirmation and customer feedback. \\
O8 & Keep customers and technicians informed through in-app notifications, email, reminders and a sanitized public tracking page. \\
O9 & Provide a technician wallet with a 90/10 settlement and controlled bKash withdrawals. \\
O10 & Give admins operational oversight: attention lists, reports, catalogue management and moderation. \\
O11 & Validate the system through automated server tests, browser tests, static checks and continuous integration, and deploy it. \\
O12 & Protect customers after the repair with a warranty, and make sure no request can get stuck in any state. \\
\bottomrule
\end{tabularx}
\end{table}

\section{Scope}
The finished system covers the following: registration and email verification, technician applications with structured expertise, a versioned service catalogue, AI-assisted request creation, damage-photo upload and pickup scheduling. It also covers an open jobs board with offers, booking with an inspection fee and refunds, inspection, quotation and revision, Stripe payment, repair progress with completion evidence, receipt confirmation, and reviews and private reports. On top of that there are in-app and email notifications, public tracking, a technician wallet with bKash withdrawals, an AI customer-support chat, admin reports with an AI summary, offer ranking, a 30-day repair warranty, deadlines and admin alerts for stalled work, a record of the device's condition at pickup, and a production deployment.

The project does not try to do everything a commercial field-service company would. Automatic mobile-wallet payouts, cash payment, spare-parts inventory, live GPS tracking of technicians, native mobile apps and accounting for several branches are outside the current scope. Chapter~\ref{ch:limitations} discusses them as future work.

\section{Individual Contribution}
Sarabo is a solo project, and I, \studentname, did all of the work on it. This covered requirements analysis, data modelling, user-interface design, frontend and backend development, the workflow and financial logic, integration of the external services, AI prompt and safety design, testing, deployment, debugging and writing this report.

\section{Report Organization}
Chapter~2 gives the technical background and looks at related systems. Chapter~3 sets out the stakeholders and requirements. Chapter~4 presents the architecture, the data model and the state model. Chapter~5 describes the implementation and how development was carried out. Chapter~6 explains the core workflows and business rules. Chapter~7 walks through the running system, with screenshots of every main step for each role. Chapter~8 covers the AI-assisted features and how they are kept safe. Chapter~9 discusses security, reliability and data integrity, and Chapter~10 presents testing and evaluation. Chapter~11 lists limitations and future work, and Chapter~12 concludes the report. The appendices summarise the API, the data model and the financial rules, give build notes, and record one notable incident from development.

\end{document}
